\documentclass[10pt,a4paper]{amsart}
\usepackage[margin=19mm]{geometry}
\usepackage{amsmath,amssymb,mathtools}
\usepackage[scr=rsfs,cal=euler]{mathalfa}
\usepackage{xfrac}
\usepackage[unicode,hidelinks,bookmarks=false]{hyperref}
\newcommand{\ie}{i.e.\ }
\newcommand{\cf}{cf.\ }
\newcommand{\resp}{respectively\ }
\newcommand{\Subsection}{\subsection}
\providecommand{\nobreakdash}{\nobreak\hbox{-}}
\setlength{\emergencystretch}{2em}
\multlinegap0pt
\usepackage{xcolor}
\usepackage{amssymb}
\usepackage{bm}
\usepackage{tikz}
\usepackage{mathtools}
\newtheorem{prop}{Proposition}
\title{A conjecture on the Crazy Knight's Tour Problem}
\author{Lorenzo Mella, Anita Pasotti}
\date{Source: arXiv:2311.09054v2 (CC BY 4.0)}
\begin{document}
\maketitle

\noindent\emph{Source and licence.} A conjecture on the Crazy Knight's Tour Problem. Version arXiv:2311.09054v2 (CC BY 4.0), distributed by arXiv under the Creative Commons Attribution 4.0 International licence, http://creativecommons.org/licenses/by/4.0/, as recorded in the arXiv licence metadata for this exact version and shown on the abstract page https://arxiv.org/abs/2311.09054v2. Full licence notice: \texttt{licenses/arxiv-2311.09054-CC-BY-4.0.txt}.\par\smallskip

\noindent\textit{Editorial scope.} Counted unit: the article's prop thm:main\_eqv (source0.tex lines 408--416). The article's complete original proof of it is delivered with it (source0.tex lines 417--487, one \texttt{proof} environment of 71 lines). No mathematics is written, repaired or re-derived: the statement and the proof are the article's own lines, in the article's own notation, and no retained line is substituted or re-flowed.  Retained uncounted as the source-local context the proof needs: lines 370--372; lines 373--375.  Nothing was omitted from any retained span, so the package carries no omission record.  No retained line contains a citation, so the wrapper carries no bibliography and no bibliography entry.  No retained line contains a figure environment, an image-inclusion command or drawing code, so no image is delivered and the package declares no asset dependency.  Editorial formatting only, in addition: an AMS \texttt{amsart} preamble that restates the article's own declarations and macros used by the retained lines, namely the environments \texttt{prop} on separate counters, together with the standard packages those lines need.  The retained environments print in this document as Proposition 1 in order of appearance, whereas the article prints the same items with the numbering of its own full document.  The complete native source of the article as distributed in the arXiv e-print for this exact version is bundled unchanged as \texttt{sources/150039/source0.tex}.
\begin{equation}\label{eq:cond1}
    j_1 < j_2 < \dotsc < j_{t+1} 
\end{equation}

\begin{equation} \label{eq:cond2}
\text{$E = \left(\bigcup_{a=1}^{t+1} I_{j_a}\right)  \setminus [j_{t+1}g-f+1, j_{t+1}g]$  for some integer $f \in [1,g-1]$,}
\end{equation}

\begin{prop} \label{thm:main_eqv}
Let $A$ be a cyclically $k$-diagonal array of size $n$, with $n>k \geq 3$ and let $g= \gcd(n,k-1)$. Let $E$ be a subset of $[1,n]$ satisfying Conditions (\ref{eq:cond1}) and (\ref{eq:cond2}) and set $h=k-1-|E|$. Moreover let $R = (r_1,\dotsc, r_n) \in \{-1,1\}^n$, where $r_i = -1$ if and only if $i \in E$, and let $C = (1,\dotsc,1)$. Then $R$ and $C$ are a solution to $P(A)$ if and only if the $h$-knight is a solution to $T(G)$, where $G:=G_f$ is
\begin{equation} \label{eq:procedureD}
\begin{aligned}
&\text{the\ $(t+1)\times (h+g)$\ array\ whose\ empty\ cells}\\
&\text{are\ exactly\ the\ last\ $f$\ cells\ of\ the\ last\ row.}
\end{aligned}
\end{equation}
\end{prop}

\begin{proof}
We divide the proof of the statement in two steps:  firstly we show  that there is a natural bijection $\phi$ between the set $E$ and the filled cells of a suitable subarray $B$ of $G$; then, we show that for every $e \in E$ it holds:
\begin{equation}\label{eq:equiv_knight_h}
S_E(e) = \phi^{-1}N_h^\alpha\phi(e),
\end{equation}
where $\alpha\geq 1$ is the minimum integer such that $N_h^\alpha(e) \in B$. From Equation (\ref{eq:equiv_knight_h}) it is then clear that the Crazy Knight's move function restricted to the set $E$ is equivalent to an $h$-knight over the array $G$, thus proving the necessary and sufficient condition of the statement.

\paragraph{\textbf{Step 1.}} From Condition (\ref{eq:cond2}) we know that $E$ is the union of $t$ intervals $I_{j_1}, \dotsc, I_{j_t}$ of size $g$, and  a  subinterval $I'_{j_{t+1}} $ of $I_{j_{t+1}}$ of size $g-f$. Let $B $ be the subarray of $G$ containing its last $g$ columns. Hence the empty cells of $B$ are exactly the last $f$ cells of its last row. Note that $|F(B)|=g(t+1)-f=gt+(g-f)=|E|.$

For every $I_{j_a} = \{e_{a,1}, \dotsc, e_{a,g}\}$ with $a \in [1,t]$ and for $I'_{j_{t+1}} = \{e_{t+1,1},\dotsc, e_{t+1,g-f} \}$, where $e_{y,x_1}<e_{y,x_2} \Leftrightarrow x_1<x_2$ for each $y \in [1,t+1]$, let $\phi$ be the map that assigns $e_{u,v}$ to the cell $(u,v)$ of $B$, which is filled from the definition of $B$. Moreover, since the intervals $I_{j_1},\dotsc,I_{j_{t}},I_{j_{t+1}}'$ are pairwise disjoint and $|F(B)| = |E|$, we have that $\phi$ is a bijection. As a final remark, notice that if we endow the indexes of the filled cells of the array $B$ with the following partial ordering:
\[
(u,v) \prec (u',v') \Leftrightarrow u<u',
\]
 the map $\phi$ then preserves the natural ordering of the set $\{j_1, \dotsc, j_{t+1}\}$.

\paragraph{\textbf{Step 2.}} 
By Condition (\ref{eq:cond1}) the sequence $(j_1,\dotsc, j_{t+1})$ is  ordered with respect to both the natural ordering of the integers and  the action of the map $\Omega$. Since $\phi$ is order-preserving, for any $e \in E$ it holds:
\[
	\phi^{-1} s_C \phi (e) = \Theta^\gamma(e).
\]
Thus, recalling the definition of $S_E$ and $N_h$, we note that Equation (\ref{eq:equiv_knight_h}) is equivalent to:
\[
\begin{aligned}
S_E(e) = \phi^{-1}N_h^\alpha\phi(e) &\Leftrightarrow  \Theta^\gamma(\Omega(e))  =  \phi^{-1}s_c s_r^h N_h^{\alpha-1}\phi(e) \\
& \Leftrightarrow \Theta^\gamma(\Omega(e))  =  \phi^{-1}s_c \phi \phi^{-1} s_r^h N_h^{\alpha-1}\phi(e).
\end{aligned}
\]
Hence, in order to prove Equation (\ref{eq:equiv_knight_h}) it is sufficient to verify that:
\begin{equation} \label{eq:simplified}
\phi^{-1} s_r^hN_h^{\alpha-1} \phi (e) = \Omega(e).
\end{equation}

Now, let $G$ be the $(t+1)\times (h+g)$ array defined in (\ref{eq:procedureD}).
Since $h=k-1-|E|$ and $|E|=g(t+1)-f$, we have $h = pg+f$ for some integer $p\geq0$, and  we prove Equation (\ref{eq:simplified}) by induction on $p$.
Assume first that $p=0$, hence $h=f$ and $G$ is a $(t+1)\times (f+g)$ array. Let $e_j \in E$ and $\phi(e_j) = (u,v)$, we distinguish between two cases:
\begin{itemize}
	\item if $u \neq t+1$ and $v\leq g-f$, or if $u = t+1$ and  $v\leq g-2f$, then
	\[
			s_R^f N_f^0\phi(e_j)  = s_R^f\phi(e_j) \in B,
	\]
	hence  $s_R^f\phi(e_j) = \phi(e_{j+f}) = \phi\Omega(e_{j})$;
	\item in the other cases, it is easy to see that:
	\[
s_R^f N_f^1\phi(e_j) \in B,
	\]
	since $s_R^f N_f^1\phi(e_j)$ is the $ (u+2f,v+1)$-th cell in the array $G$, this corresponds to the $(u+f,v+1)$-th cell in the array $B$, that is exactly $\phi\Omega(e_j)$.
\end{itemize}

Assume now that Equation (\ref{eq:simplified}) holds for $h = pg+f$; we prove that this implies that it holds for $h'= (p+1)g+f$. Let $G$ and $G'$ be the two arrays relative to $h$ and $h'$, respectively. It can be easily seen that $G'$ is composed by the concatenation of a completely filled $(t+1)\times g$ array with $G$, hence we refer to $G$ as being contained inside $G'$. Then, let $\Omega$ and $\Omega'$ be the functions on $E$ relative to $h$ and $h'$, respectively.


 
Assume first that $e \in E$ is such that $s_c \phi\Omega (e)  =\phi \Omega'(e)$.  By inductive hypothesis, on $G$ we have $ s_r^hN_h^{\alpha-1} \phi (e) = \phi\Omega(e)$. 
We point out that in $G$  (resp. $G'$) an $h$-knight (resp. $h'$-knight) move is equivalent to a $(-g)$-knight move for all rows, except for the last one where it corresponds to a $(-g+f)$-knight move.
Let $(u,v) = N_h^{\alpha-1} \phi (e)$ in $G$, that is also $N_{h'}^{\alpha-1} \phi (e)$ in $G'$ by previous remark: it is easy to see that, since $s_r^h(u,v)$ in $G$ belongs to the array $B$, the move $s_r^{h'}(u,v)$ is contained in $G' \setminus G$, hence we apply $N_{h'}$ to $(u,v)$. The resulting cell belongs to $G' \setminus G$, and it follows that $s_r^{h'}N_{h'}(u,v)$ is contained in $B$. Hence, in $G'$, it holds:
\[
s_r^{h'}N_{h'}(u,v) = s_r^{h'}N_{h'}N_{h'}^{\alpha-1}\phi(e) =  s_r^{h'}N_{h'}^{\alpha}\phi(e)  \in B.
\] 
Comparing this formula to the one  for $G$, we have applied  an $h'$-knight move once more. A simple counting argument shows that $s_r^{h'}N_{h'}^{\alpha}\phi(e)$ in $G'$ belongs to the same column of  $ s_r^hN_h^{\alpha-1} \phi (e) = \phi\Omega(e)$ in $G$, and as it is contained in the successive row it follows that:
 \[
 s_r^{h'}N_{h'}^{\alpha}\phi(e) = s_c \phi\Omega (e)  =\phi \Omega'(e).
 \]

Let now $e \in E$ be such that $s_c \phi\Omega (e)  \neq \phi \Omega'(e)$. Notice that this happens if and only if  $\Omega (e) = e_u$, where $u \in[|E|-g+1,|E|]$. Also in this case, let $(u,v) =N_h^{\alpha-1}	\phi(e)$ in $G$, that is also $N_{h'}^{\alpha-1}\phi(e)$ in $G'$. It easily follows that $v \in \{t,t+1\}$, and that $s_r^{h'}(u,v)$ does not belong to $B$, hence we apply the $h'$-knight move function.

If $v = t+1$, then $N_{h'}(u,v) \in G' \setminus G$, in particular, it belongs to the first row of $G'$, hence $s_r^{h'}N_{h'}(u,v) = (x_1,y_1) \in B$. 
If $v = t$, then $N_{h'}(u,v) \in G' \setminus G$; in particular, it belongs to the last row of $G'$. Since $s_c \phi\Omega (e)  \neq \phi \Omega'e)$, it follows that $s_r^{h'}N_{h'}(u,v) \not\in B$, hence we apply the $h'$-knight move function another time, thus $s_r^{h'} N^2_{h'}(u,v) = (x_2,y_2)$ belongs to $B$. Let $\eta $ be the map that associates to each $e \in E $  such that $s_c \phi\Omega (e)  \neq \phi \Omega'(e)$ its corresponding cell, that is either $(x_1,y_1)$ or $(x_2,y_2)$.
It is easy to see that $x_1 = x_2 = 1$, and as $|[|E|-g+1,|E|]| = g$ we have that $\eta$ is a bijection between the elements of $[|E|-g+1,|E|]$ and the first row of $B$. To conclude, it is sufficient to check that $\eta$ is order-preserving, where the ordering of the cells of the row is intended as the natural one, i.e. $(1,y_1) \prec (1,y_2) \Leftrightarrow y_1 <y_2$. 
Let $(u_\gamma,v_\gamma) = N_{h'}^{\alpha_\gamma-1}\phi(e_\gamma)$ for $\gamma \in \{1,2\}$, where $e_1 < e_2$. Clearly, if $u_1 = u_2$, then  the order is preserved. If $u_1 = t$ and $u_2 = t+1$, then $N_h'(u_1,v_1) < (u_2,v_2)$: in fact, $N_h'(u_1,v_1)$ belongs to $G' \setminus G$, while $(u_2,v_2)$ belongs to $G$.
Thus, $\eta$ is order-preserving, concluding the proof.
\end{proof}

\end{document}
